\subsection{CONVERGENCE CRITERIA FOR SERIES WITH POSITIVE TERMS}

In this section by a series with positive terms we shall understand (by an abuse of language) a series \((u_{n})\) of real terms such that \(u_{n} \geqslant 0\) starting from some particular value of \(n\) . Everything we shall say about such series will extend immediately by a change of sign to series all of whose terms are \(\leqslant 0\) starting from some particular value of \(n\) . We have seen (II, p.~64, Example 3) that to every sequence \((\mathbf{u}_n)_{n\geqslant 1}\) of points in a normed space \(E\) one can associate a step function \(\mathbf{u}\) defined on \([1, +\infty[\) by the conditions \(\mathbf{u}(x) = \mathbf{u}_n\) for \(n\leqslant x < n + 1\) : then the series \((\mathbf{u}_n)\) converges if and only if the integral \(\int_1^{+\infty}\mathbf{u}(t)dt\) converges.

Let \((u_{n})\) and \((v_{n})\) be two series with positive terms, and u and v the associated step functions: the relation \(u_{n} \leqslant v_{n}\) for \(n \geqslant n_{0}\) is equivalent to \(u(x) \leqslant v(x)\) for \(x \geqslant n_{0}\) . Thus each of the relations \(u_{n} \preccurlyeq v_{n}, u_{n} \prec v_{n}, u_{n} \sim v_{n}\) is equivalent to \(u(x) \preccurlyeq v(x), u(x) \prec v(x)\) and \(u(x) \sim v(x)\) respectively; this remark allows us to translate propositions 1 (V, p.~228) and 6 (V, p.~230) as follows:

PROPOSITION 1. Let \((u_{n})\) and \((v_{n})\) be two series with positive terms. If \(u_{n} \preccurlyeq v_{n}\) and if the series \((v_{n})\) converges, then \((u_{n})\) converges; if \(u_{n} \succcurlyeq v_{n}\) and if \(\sum_{n=1}^{\infty} v_{n} = +\infty\) , then \(\sum_{n=1}^{\infty} u_{n} = +\infty\) .

PROPOSITION 2. Let \((u_{n})\) and \((v_{n})\) be two series with positive terms:

\(1^{\circ}\) If the series \((v_{n})\) converges then the relation \(u_{n} \ll v_{n}\) (resp. \(u_{n} \sim v_{n}\) ) implies \(\sum_{p=n}^{\infty} u_{p} \ll \sum_{p=n}^{\infty} v_{p}\) (resp. \(\sum_{p=n}^{\infty} u_{p} \sim \sum_{p=n}^{\infty} v_{p}\) ).

\(2^{\circ}\) If \(\sum_{n=1}^{\infty} v_n = +\infty\) then the relation \(u_n \ll v_n\) (resp. \(u_n \sim v_n\) ) implies \(\sum_{p=1}^{n} u_p \ll \sum_{p=1}^{n} v_p\) (resp. \(\sum_{p=1}^{n} u_p \sim \sum_{p=1}^{n} v_p\) ).

One obtains convenient convergence criteria by taking for the comparison series \((v_{n})\) in prop. 1 a series whose terms are of the form \(v_{n} = f(n)\) , where \(f\) is a function \(\geqslant 0\) defined for every real number \(x > x_0\) and decreasing on the interval \([x_0, +\infty[\) ; indeed:

PROPOSITION 3 (Cauchy-Maclaurin criterion). If \(f\) is a function \(\geqslant 0\) and decreasing on \([x_0, +\infty[\) , then the series with general term \(v_n = f(n)\) converges if and only if the integral \(\int_{x_0}^{+\infty} f(t) dt\) converges.

To prove this it is sufficient to note that if v is the step function associated with the series \((v_{n})\) then \(v(x+1) \leqslant f(x) \leqslant v(x)\) for every \(x \geqslant x_{0}\) since f is decreasing; the proposition is thus a consequence of the comparison principle (II, p.~66, prop. 3).

Since the functions which feature in the logarithmic convergence criteria for integrals (V, p.~229, prop. 2, 3 and 4) are decreasing on an interval \([x_{0}, +\infty[\) , applying prop. 1 and 3 of V, p.~237, gives the following criteria:

PROPOSITION 4 (``logarithmic criterion of order 0''). Let \((u_{n})\) be a series with positive terms; if \(u_{n} \preccurlyeq n^{\mu}\) for some \(\mu < -1\) then the series \((u_{n})\) converges; if \(u_{n} \succcurlyeq n^{\mu}\) for some \(\mu \geqslant -1\) then the series \((u_{n})\) has an infinite sum.

PROPOSITION 5 (``logarithmic criterion of order p''). Let \((u_{n})\) be a series with positive terms. If \(u_{n} \preccurlyeq \frac{1}{nl_{1}(n)l_{2}(n)\ldots l_{p-1}(n)(l_{p}(n))^{\mu}}\) for some \(\mu > 1\) then the series \((u_{n})\) converges; if \(u_{n} \succcurlyeq \frac{1}{nl_{1}(n)l_{2}(n)\ldots l_{p-1}(n)(l_{p}(n))^{\mu}}\) for some \(\mu \leqslant 1\) then the series \((u_{n})\) has an infinite sum.

If \(0 \leqslant q < 1\) one has \(q^n \preccurlyeq n^{-\mu}\) for any \(\mu > 0\) ; applying the logarithmic criterion of order 0 again proves the convergence of the geometric series \(\sum_{n=0}^{\infty} q^n\) for \(|q| < 1\) (Gen.~Top., IV, p.~364). If one applies prop. 1 with \(v_n = q^n\) one obtains a criterion which may be put in the following form (``Cauchy's criterion''): Let \((u_n)\) be a series with positive terms; if \(\limsup_{n \to \infty} (u_n)^{1/n} < 1\) then the series \((u_n)\) converges; if \(\limsup_{n \to \infty} (u_n)^{1/n} > 1\) then the series \((u_n)\) has infinite sum. Indeed, if \(\limsup_{n \to \infty} (u_n)^{1/n} = a < 1\) then \(u_n \preccurlyeq q^n\) for every \(q\) such that \(a < q < 1\) . If, on the other hand, \(\limsup_{n \to \infty} (u_n)^{1/n} = a > 1\) then \(u_n \geqslant q^n > 1\) for infinitely many values of \(n\) , for any \(q\) such that \(1 < q < a\) ; since \(u_n\) does not tend to 0 one has \(\sum_{n=1}^{\infty} u_n = +\infty\) .

This criterion is very useful in the theory of entire series, which we shall study later; but it even so it does not permit one to decide the convergence of the series \((1/n^{\alpha})\) , in other words, its field of application is much more restricted than that of the logarithmic criteria.

